%----------------------------------------------------------------------------------------
% METADATA BLOCK (YAML)  — see inbox/README.md
%----------------------------------------------------------------------------------------
% uid: ""
% school: "lycee"
% level: "3eme"
% section: "Economie et Gestion"
% chapter: "Statistiques"
% title: "Quelle machine est la plus régulière ?"
% difficulty: "Facile"
% tags: ["statistiques", "variance", "ecart-type", "coefficient-de-variation"]
% provenance: "own"
% batch: "2026-11-19-3eco-statistiques-n5"
%----------------------------------------------------------------------------------------

\newcommand{\R}{\mathbb{R}}
\newcommand{\vect}[1]{\overrightarrow{#1}}

% ---STATEMENT---
Une usine contrôle $20$ paquets remplis par chacune de ses deux machines. Masses en grammes :
\begin{center}
\begin{tabular}{|l|c|c|c|c|c|}\hline
Machine A (paquets de $250$ g) & $246$ & $248$ & $250$ & $252$ & $254$ \\\hline
Effectif & $2$ & $4$ & $8$ & $4$ & $2$ \\\hline
\end{tabular}
\end{center}

\begin{center}
\begin{tabular}{|l|c|c|c|c|c|}\hline
Machine B (paquets de $1\,000$ g) & $994$ & $997$ & $1\,000$ & $1\,003$ & $1\,006$ \\\hline
Effectif & $1$ & $4$ & $10$ & $4$ & $1$ \\\hline
\end{tabular}
\end{center}
\begin{enumerate}
  \item Calculer la masse moyenne de chaque série.
  \item Calculer la variance et l'écart type de chaque série (au centième).
  \item Calculer le coefficient de variation $\dfrac{\sigma}{\overline x}$ de chaque série (en pourcentage, au centième).
  \item Quelle machine est la plus régulière ? Justifier.
\end{enumerate}

% ---SOLUTION---
\begin{enumerate}[itemsep=14pt]
  \item Les deux tableaux sont symétriques autour de leur valeur centrale :
  \[\overline x_A=\dfrac{2\times246+4\times248+8\times250+4\times252+2\times254}{20}=\dfrac{5\,000}{20}=250\text{ g}\]
  \[\overline x_B=\dfrac{994+4\times997+10\times1\,000+4\times1\,003+1\,006}{20}=\dfrac{20\,000}{20}=1\,000\text{ g}\]

  \item
  \begin{prop}[Variance avec les écarts]
  \ib{$V=\dfrac{n_1(x_1-\overline x)^2+\dots+n_p(x_p-\overline x)^2}{N}$} : elle donne le même résultat que « moyenne des carrés $-\ \overline x^{\,2}$ », avec des nombres plus petits.
  \end{prop}
  Machine A, écarts à $250$ : $-4$, $-2$, $0$, $2$, $4$.
  \[V_A=\dfrac{2\times16+4\times4+0+4\times4+2\times16}{20}=\dfrac{96}{20}=4{,}8\qquad\sigma_A=\sqrt{4{,}8}\approx2{,}19\text{ g}\]
  Machine B, écarts à $1\,000$ : $-6$, $-3$, $0$, $3$, $6$.
  \[V_B=\dfrac{36+4\times9+0+4\times9+36}{20}=\dfrac{144}{20}=7{,}2\qquad\sigma_B=\sqrt{7{,}2}\approx2{,}68\text{ g}\]

  \item
  \begin{prop}[Coefficient de variation]
  \ib{$CV=\dfrac{\sigma}{\overline x}$} : il compare la dispersion de séries de moyennes différentes.
  \end{prop}
  \[CV_A=\dfrac{2{,}19}{250}\approx0{,}88\,\%\qquad CV_B=\dfrac{2{,}68}{1\,000}\approx0{,}27\,\%\]

  \item $\sigma_B>\sigma_A$, mais les moyennes sont très différentes : il faut comparer les coefficients de variation.

  $CV_B<CV_A$ : \textbf{la machine B est la plus régulière} (relativement à la masse qu'elle doit remplir).
\end{enumerate}
\textbf{Piège :} conclure avec les écarts types seuls ; une erreur de $2{,}7$ g sur $1$ kg est plus petite, en proportion, qu'une erreur de $2{,}2$ g sur $250$ g.
